% Generated draft: requires mathematical and source review.
\documentclass[11pt]{article}
\usepackage[margin=25mm]{geometry}
\usepackage{fontspec}
\setmainfont{Latin Modern Roman}
\usepackage{amsmath,amssymb,mathtools,unicode-math}
\setmathfont{Latin Modern Math}
\usepackage{enumitem,needspace,xurl,hyperref}
\setlength{\parindent}{0pt}
\setlength{\parskip}{5pt}
\setlength{\emergencystretch}{3em}
\newcommand{\N}{\mathbb{N}}
\newcommand{\Z}{\mathbb{Z}}
\newcommand{\Q}{\mathbb{Q}}
\newcommand{\R}{\mathbb{R}}
\newcommand{\C}{\mathbb{C}}
\newcommand{\F}{\mathbb{F}}
\newcommand{\A}{\mathbb{A}}
\newcommand{\PP}{\mathbb{P}}
\newcommand{\E}{\mathbb{E}}
\newcommand{\eps}{\varepsilon}
\newcommand{\dd}{\,\mathrm{d}}
\begin{document}
\section*{UnsolvedMath ID: 801100}
\textbf{Source problem code:} catalogue entry
\subsection{Definitions and mathematical statement}
Let \(\operatorname{Ult}(\N)\) denote the set of ultrafilters on \(\N\). An ultrafilter is a family \(\mathcal U\subseteq\mathcal P(\N)\) such that \(\N\in\mathcal U\), \(\varnothing\notin\mathcal U\), it is closed under finite intersections and passage to supersets, and for every \(S\subseteq\N\), exactly one of \(S\) and \(\N\setminus S\) belongs to \(\mathcal U\).

For structures \(A,B\) in a fixed first-order language \(L\), elementary equivalence, denoted \(A\equiv_L B\), means that every \(L\)-sentence has the same interpretation in \(A\) and \(B\). In continuous logic this means equality of the values of every sentence; in classical logic it means agreement on the truth of every sentence.

\textbf{Requested classification, conditional on unresolved notation.} The source asks to classify the elementary equivalence classes of
\[
\bigl\{\mathbb M^{\mathcal U}:\mathcal U\in\operatorname{Ult}(\N)\bigr\},
\]
and asks the analogous question for
\[
\bigl\{\mathbb M_{\mathcal U}:\mathcal U\in\operatorname{Ult}(\N)\bigr\}.
\]
Once these constructions and their common language \(L\) are specified, the respective equivalence relations on the indexing ultrafilters are
\[
\mathcal U\sim^{\mathrm{sup}}_L\mathcal V
\quad\Longleftrightarrow\quad
\mathbb M^{\mathcal U}\equiv_L\mathbb M^{\mathcal V},
\qquad
\mathcal U\sim^{\mathrm{sub}}_L\mathcal V
\quad\Longleftrightarrow\quad
\mathbb M_{\mathcal U}\equiv_L\mathbb M_{\mathcal V}.
\]
The task is to classify the two resulting families up to their respective elementary equivalence relations.

\textbf{Unresolved qualifications.} The supplied record does not define \(\mathbb M\), \(\mathbb M^{\mathcal U}\), or \(\mathbb M_{\mathcal U}\), nor does it specify the logical language. The title supplies a context of ultraproduct \(\mathrm{II}_1\) factors, but does not determine these constructions. The superscript and subscript notation must therefore remain distinct and uninterpreted here. The literal phrase ``ultrafilters on \(\N\)'' includes principal ultrafilters; whether the intended range is restricted to nonprincipal ultrafilters is not stated. These missing specifications prevent an unconditional precise formulation of the intended classification problem.
\subsection{Short English statement}
The source asks for a classification, up to elementary equivalence, of each of the two families \(\mathbb M^{\mathcal U}\) and \(\mathbb M_{\mathcal U}\) as \(\mathcal U\) ranges over ultrafilters on \(\N\). The constructions, logical language, and any intended restriction to nonprincipal ultrafilters require clarification before the mathematical target is fully determined.
\subsection{Sources}
\begin{itemize}
\item Local UnsolvedMath registry, permanent ID \(801100\), exact source problem code \texttt{catalogue entry}, title ``Ultraproduct II\_1 Factors --- Problem 1.'' The supplied record attributes the problem to Elayavalli, with proposed year 2025, and labels its status as open. This status is supplied registry information, not an independently verified mathematical assessment.
\item Supplied source attribution: Elayavalli, \emph{50 Open Problems: Ultraproduct II\_1 Factors} (2025), Problem 1. Supplied source URL: \url{https://arxiv.org/abs/2511.20377}. The record reports access on 2026-10-16 and extraction by \texttt{arxiv\_latex}; these are supplied provenance details, not a claim of independent source access or verification.
\end{itemize}
\end{document}
